\section{Related Work}
\label{sec:related}

% \paragraph{3D Generation with 2D Generative Models.}
% SDS-style text-to-3D leverages frozen 2D diffusion priors to optimize differentiable 3D representations (\emph{e.g.}, NeRFs), introduced by DreamFusion~\cite{dreamfusion} and SJC~\cite{sjc}, with subsequent variants improving resolution, disentangling geometry/appearance, mitigating saturation/diversity issues, and adding multi-view/geometry regularizers; recent hybrids with 3D Gaussian Splatting or 3D diffusion further improve speed and fidelity~\cite{consistent3d,magic3d,fantasia3d,prolificdreamer,li2024connecting,kwak2024geometry,dreamcraft3d,gaussiandreamer,dreamgaussian, liang2024luciddreamer, li2023instant3d}.
% % A parallel thread adapts SDS guidance to single-image 3D reconstruction—either by direct optimization or via view-conditioned diffusion followed by reconstruction—with improved coherence and runtime in later systems~\cite{realfusion,neurallift360,makeit3d,liu2023zero,weng2023consistent123,liu2024one,liu2023one,wonder3d,syncdreamer,xu2024instantmesh,li2023instant3d}.
% % Despite steady progress, reliance on 2D priors often leaves multi-view inconsistencies, motivating native 3D generative models. 

% A parallel line of work extends SDS-style guidance to single-image 3D reconstruction, either through direct optimization of a 3D representation or through view-conditioned diffusion followed by 3D reconstruction. These methods have progressively improved cross-view coherence and runtime efficiency~\cite{realfusion,neurallift360,makeit3d,liu2023zero,weng2023consistent123,liu2024one,liu2023one,wonder3d,syncdreamer,xu2024instantmesh,li2023instant3d}. Despite this steady progress, methods built on 2D priors still often suffer from limited 3D consistency, which has motivated growing interest in native 3D generative models.

% Among the works most closely related to ours are TICD~\cite{ticd} and FlexGen~\cite{flexgen}. TICD augments a text-to-3D SDS pipeline with an additional image diffusion prior, and FlexGen takes both text and image prompts as input but focuses on 2D multi-view image generation rather than native 3D synthesis. In contrast, we study native 3D generation jointly conditioned on text and image inputs.

\paragraph{3D Generation with 2D Generative Models.}
SDS-based text-to-3D methods optimize differentiable 3D representations (\emph{e.g.}, NeRFs~\cite{nerf}) using frozen 2D diffusion priors, as pioneered by DreamFusion~\cite{dreamfusion} and SJC~\cite{sjc}. Later works improve resolution, geometry--appearance disentanglement, optimization stability, and multi-view consistency, while recent hybrids with 3D Gaussian Splatting or native 3D diffusion further boost efficiency and fidelity~\cite{consistent3d,magic3d,fantasia3d,prolificdreamer,li2024connecting,kwak2024geometry,dreamcraft3d,gaussiandreamer,dreamgaussian,liang2024luciddreamer,li2023instant3d}. A parallel line studies image-conditioned 3D generation and reconstruction with 2D generative priors, either by directly optimizing 3D representations or by first generating consistent multi-view images and then reconstructing 3D content~\cite{realfusion,neurallift360,makeit3d,liu2023zero,weng2023consistent123,liu2024one,liu2023one,wonder3d,syncdreamer,xu2024instantmesh,li2023instant3d}. Despite strong progress, reliance on 2D priors often limits 3D consistency, motivating native 3D generative models and cross-modal latent modeling~\cite{michelangelo}.

Among prior works, TICD~\cite{ticd} and FlexGen~\cite{flexgen} are most related to ours. TICD augments SDS-based text-to-3D with an image diffusion prior, while FlexGen jointly conditions on text and image but focuses on 2D multi-view generation rather than native 3D synthesis. In contrast, we study \textbf{native} 3D generation under joint text-image conditioning.

\paragraph{Native 3D Generative Models.}
Unlike SDS-based pipelines that rely on 2D generators at test time, native 3D generative models operate directly on 3D representations such as point clouds, meshes, voxels, 3D Gaussians~\cite{3dgs,tang2024lgm}, and neural fields~\cite{rf, nerf, mipnerf, lan2024ln3diff}. Early works~\cite{nichol2022point, vahdat2022lion} introduce fast point-cloud synthesis via latent diffusion. Later models such as 3DShape2VecSet~\cite{zhang20233dshape2vecset} improve geometry-aware latents for diffusion training. Moreover, scalable voxel-based methods~\cite{ren2024xcube}, octree-based models~\cite{xiong2025octfusion} and 3DGS-based models~\cite{yang2024atlas} improve generation resolution and efficiency. 

More recent studies scale data and model capacity for high-quality asset synthesis. Some focus on geometry-only generation for detailed meshes~\cite{chen2025ultra3d,direct3d,direct3ds2,li2025sparc3d,ye2025hi3dgen}, while others target fully textured 3D assets~\cite{clay,3dtopiaxl,hunyuan3d2025hunyuan3d2.1,lai2025hunyuan3d2.5,step1x,zhou2025few}. Among the latter, TRELLIS~\cite{trellis} introduces a sparse, structured latent representation that enables high-fidelity 3D generation via a two-stage pipeline, and UniLat3D~\cite{unilat} provides a unified-latent variant with a simplified sampling process. Together, these advances establish native 3D diffusion models as strong alternatives to 2D diffusion-based pipelines. 
However, most existing approaches assume single-modality conditioning (either image or text), which limits the flexibility and expressiveness of user instructions. In this work, we study this limitation and introduce the task of text–image conditioned 3D generation.

% \paragraph{Multimodal-Conditioned Generation.}
% Multimodal conditioning has become a powerful recipe across text, image, and video generation. In text generation, multimodal large language models (MLLMs) that accept textual, visual~\cite{llava,qwenvl,instructblip,internvl} and even 3D inputs (\emph{e.g.}, RGB-D or point clouds)~\cite{3dllm,pointllm,ll3da} demonstrate grounded reasoning and control, advancing visual understanding and robotics~\cite{rt2,palme,pi0}. For image synthesis, explicit visual controls (edges, depth, pose, layout) injected into text-to-image diffusion~\cite{controlnet,t2i,ipadapter} substantially improve alignment with user intent. Joint text–image conditioning~\cite{qwenimage,bagel,instructp2p} also enables instruction-following and layout-grounded image editing in a single forward pass. In video generation, conditioning on a reference image (or first frame) while guiding subsequent frames with text~\cite{makeitmove, wan} is an emerging strategy for temporally consistent, controllable outputs. These successes motivate our exploration of whether joint text-image conditioning can endow 3D generators with complementary strengths beyond single-modality conditioning.

\paragraph{Multimodal-Conditioned Generation.}
Multimodal conditioning has proven effective in text, image, and video generation. Multimodal LLMs that ingest both language and visual (or even 3D) inputs~\cite{llava,qwenvl,instructblip,internvl,3dllm,pointllm,ll3da} enable grounded reasoning and control, while text-to-image diffusion benefits from additional visual controls (edges, depth, pose, layout) and joint text-image inputs~\cite{controlnet,t2i,ipadapter,qwenimage,bagel,instructp2p} for better alignment with user intent. In video generation, combining a reference frame with textual guidance~\cite{makeitmove,wan} yields more consistent and controllable outputs. These successes motivate our exploration of whether joint text-image conditioning can similarly endow 3D generators with complementary strengths beyond single-modality conditioning.